\documentclass[10pt,a4paper]{amsart}
\usepackage[margin=19mm]{geometry}
\usepackage{amsmath,amssymb,mathtools}
\usepackage[scr=rsfs,cal=euler]{mathalfa}
\usepackage{xfrac}
\usepackage[unicode,hidelinks,bookmarks=false]{hyperref}
\newcommand{\ie}{i.e.\ }
\newcommand{\cf}{cf.\ }
\newcommand{\resp}{respectively\ }
\newcommand{\Subsection}{\subsection}
\providecommand{\nobreakdash}{\nobreak\hbox{-}}
\setlength{\emergencystretch}{2em}
\multlinegap0pt
\usepackage{amssymb}
\usepackage{mathtools}
\usepackage{xcolor}
\usepackage[shortlabels]{enumitem}
\newtheorem{lemma}{Lemma}
\newtheorem{theorem}{Theorem}
\DeclareMathOperator{\E}{\mathbb{E}}
\title{MaxCut for \(\mathrm{MTP}_2\) Covariances}
\author{Gleb Smirnov}
\date{Source: arXiv:2609.10474v1 (CC BY 4.0)}
\begin{document}
\maketitle

\noindent\emph{Source and licence.} MaxCut for \(\mathrm{MTP}_2\) Covariances. Version arXiv:2609.10474v1 (CC BY 4.0), distributed by arXiv under the Creative Commons Attribution 4.0 International licence, http://creativecommons.org/licenses/by/4.0/, as recorded in the arXiv licence metadata for this exact version and shown on the abstract page https://arxiv.org/abs/2609.10474v1. Full licence notice: \texttt{licenses/arxiv-2609.10474-CC-BY-4.0.txt}.\par\smallskip

\noindent\textit{Editorial scope.} Counted unit: the article's theorem thm:quantum (source0.tex lines 426--455). The article's complete original proof of it is delivered with it (source0.tex lines 457--566, one \texttt{proof} environment of 110 lines). No mathematics is written, repaired or re-derived: the statement and the proof are the article's own lines, in the article's own notation, and no retained line is substituted or re-flowed.  Retained uncounted as the source-local context the proof needs: lines 334--358; lines 396--408.  Nothing was omitted from any retained span, so the package carries no omission record.  No retained line contains a citation, so the wrapper carries no bibliography and no bibliography entry.  No retained line contains a figure environment, an image-inclusion command or drawing code, so no image is delivered and the package declares no asset dependency.  Editorial formatting only, in addition: an AMS \texttt{amsart} preamble that restates the article's own declarations and macros used by the retained lines, namely the environments \texttt{lemma}, \texttt{theorem} on separate counters, together with the standard packages those lines need.  The retained environments print in this document as Lemma 1, Theorem 1 in order of appearance, whereas the article prints the same items with the numbering of its own full document.  The complete native source of the article as distributed in the arXiv e-print for this exact version is bundled unchanged as \texttt{sources/209921/source0.tex}.
\begin{lemma}\label{total}
Let \(\psi\) be a unit state. 
Let \((P_z)_{z \in \mathcal Z}\) be
orthogonal projections such that:
\[
\sum_zP_z=I.
\]
Set:
\[
p_z=(\psi,P_z\psi).
\]
For \(p_z>0\), let
\[
\psi_z=\frac{P_z\psi}{\sqrt{p_z}}
\]
be the state conditioned on the measurement outcome \(z\).
If \(A\) commutes with every \(P_z\), then:
\[
(\psi,A\psi)
=
\sum_{z}
p_z(\psi_z,A\psi_z),
\]
where terms \(p_z=0\) are omitted.
\end{lemma}

\begin{lemma}\label{joint}
Let \(\psi\) be a unit state. Measure the commuting observables \(Y_1,\ldots,Y_n\) simultaneously in the state \(\psi\). Since each \(Y_i\) has eigenvalues \(\pm1\), 
the outcome is a random vector
\[
S=(S_1,\ldots,S_n)\in\{-1,1\}^n.
\]
Then:
\[
\E[S_iS_j]
=
(\psi,Y_iY_j\psi).
\]
\end{lemma}

\begin{theorem}\label{thm:quantum}
Let
\[
\psi
=
\sum_{x\in\{0,1\}^n}
\psi(x)e_x
\]
be a positive log-supermodular state. Define \(L=L(\psi)\in\mathbb R^{n\times n}\) by
\[
L_{ii}=0,
\qquad
L_{ij}=-(\psi,Y_iY_j\psi)
\quad \text{for \(i \neq j\)}.
\]
Then:
\[
L_{ij} \ge 0
\]
and for each collection of nonnegative weights \(w_{ij}\),
\[
\sum_{i<j}w_{ij}L_{ij}
\le
2\,\operatorname{MaxCut}(G,w)-W,
\]
where
\[
W=\sum_{i<j}w_{ij}.
\]
\end{theorem}

\begin{proof}
Fix distinct \(i,j\in[n]\). 
Measure all qubits other than \(i,j\)
in the computational basis.
An outcome is a configuration
\[
z\in\{0,1\}^{[n]\setminus\{i,j\}}.
\]
\(z\) specifies the values of all coordinates
other than \(i\) and \(j\). Let
\[
H_z
=
\operatorname{span}\{e_{abz}:a,b\in\{0,1\}\},
\]
where \(e_{abz}\) is the computational-basis vector
with \(i\)-th coordinate \(a\),
\(j\)-th coordinate \(b\),
and remaining coordinates \(z\).
Let \(P_z\) be the orthogonal projection onto \(H_z\).
Set:
\[
p_z=(\psi,P_z\psi).
\]
For \(p_z>0\), the conditioned state is
\begin{equation}\label{psi_z}
\psi_z
=
\frac{P_z\psi}{\sqrt{p_z}}
=
\alpha_z e_{11z}
+\beta_z e_{10z}
+\gamma_z e_{01z}
+\delta_z e_{00z}.
\end{equation}
Log-supermodularity on this two-dimensional face gives:
\[
\alpha_z\delta_z
\ge
\beta_z\gamma_z.
\]
\smallskip%

On \(H_z\), the operator \(Y_iY_j\) acts by
\[
Y_iY_j e_{11z}=-e_{00z},
\quad
Y_iY_j e_{10z}=e_{01z},
\quad
Y_iY_j e_{01z}=e_{10z},
\quad
Y_iY_j e_{00z}=-e_{11z}.
\]
Hence:
\[
(\psi_z,Y_iY_j\psi_z)
=
2(\beta_z\gamma_z-\alpha_z\delta_z)
\le0.
\]
Since \(Y_iY_j\) acts only on qubits \(i,j\),
it commutes with every \(P_z\). Lemma~\ref{total} gives:
\[
-L_{ij}
=
(\psi,Y_iY_j\psi)
=
\sum_z
p_z(\psi_z,Y_iY_j\psi_z)
\le0.
\]
Let again
\[
S=(S_1,\ldots,S_n)\in\{-1,1\}^n
\]
be the simultaneous measurement of 
\(Y_1, \ldots, Y_n\). By Lemma~\ref{joint},
\[
\E[S_iS_j]
=
(\psi,Y_iY_j\psi)
=
-L_{ij},
\]
while
\[
\E[S_i^2]=1.
\]
Let \(w_{ij}\ge0\). For \(s\in\{-1,1\}^n\), write:
\[
\operatorname{Cut}_w(s)
=
\sum_{i<j}
w_{ij}\frac{1-s_is_j}{2}.
\]
Then:
\[
-\sum_{i<j}w_{ij}s_is_j
=
2\operatorname{Cut}_w(s)-W.
\]
Consequently,
\[
\sum_{i<j}w_{ij}L_{ij} = 
\E\!\left[
2\operatorname{Cut}_w(S)-W
\right] \le 2\operatorname{MaxCut}(G,w)-W,
\]
and the proof follows.
\end{proof}

\end{document}
